\section{Sujet TI n\textsuperscript{o}\,3 — Cybercafé}

\begin{center}
\setlength{\fboxsep}{8pt}
\fcolorbox{onbline}{white}{%
\begin{minipage}{0.92\linewidth}
\small\sffamily
\begin{tabular}{@{}p{0.30\linewidth}@{\hfill}p{0.34\linewidth}@{\hfill}p{0.28\linewidth}@{}}
\textbf{Office du Baccalauréat}\newline du Cameroun\newline Session 2025
&
\textbf{Série} : TI\newline \textbf{Épreuve} : Algorithmique\newline et Programmation
&
\textbf{Durée} : 03 heures\newline \textbf{Coef.} : 04
\\
\end{tabular}
\par\smallskip
{\color{onbline}\rule{\linewidth}{0.8pt}}
\par\smallskip
\centering\textbf{ÉPREUVE D'ALGORITHMIQUE ET PROGRAMMATION}
\end{minipage}}
\end{center}

\begin{remarque}
Sujet type \emph{original} OnBuch, calqué sur la structure réelle de l'épreuve
d'Algorithmique et Programmation de la série~TI (deux parties de 10~points :
programmation web puis programmation procédurale en~C). Les corrigés qui suivent
la mention « Corrigé » sont proposés par OnBuch.
\end{remarque}

\noindent\textbf{Contexte.} Le \textit{Cyber « LumièreNet »}, situé au quartier
Briqueterie à Yaoundé, loue des postes de connexion Internet aux élèves et
étudiants. Le gérant souhaite informatiser le suivi de ses sessions de connexion.
Pour chaque session, on note le \textbf{numéro de poste}, le \textbf{nom du
client}, l'\textbf{heure de début}, la \textbf{durée} (en minutes) et le
\textbf{montant} facturé (en FCFA). Le tarif appliqué est de \textbf{300~FCFA
l'heure}. On vous demande, à travers les deux parties suivantes, de l'aider à
construire l'outil web puis le module en langage~C.

% =====================================================================
\subsection*{Partie I — Programmation web \hfill (10 points)}

\textbf{Exercice 1 — Page d'accueil HTML et JavaScript.} \hfill \textbf{(4 points)}

\smallskip
Le gérant veut une page d'accueil contenant un petit \textbf{simulateur de coût}
qui, à partir d'une durée saisie en minutes, affiche le montant à payer (au tarif
de 300~FCFA l'heure), sans recharger la page.

\begin{enumerate}[label=\textbf{\arabic*-},leftmargin=2.2em]
\item Définir : \textbf{page web statique}, \textbf{balise}, \textbf{événement
\code{onclick}}. \hfill \textbf{0,75pt}
\item Donner la différence fondamentale entre un script \textbf{JavaScript} et un
script \textbf{PHP} quant au lieu de leur exécution. \hfill \textbf{0,5pt}
\item Écrire le code \textbf{HTML} d'un formulaire contenant : un champ de saisie
de durée (en minutes) d'identifiant \code{duree}, un bouton « Calculer » et une
zone de résultat d'identifiant \code{resultat}. \hfill \textbf{1,5pt}
\item Écrire la fonction \textbf{JavaScript} \code{calculerCout()} qui lit la
durée saisie, calcule le montant \emph{(montant $=$ durée $\div\,60 \times 300$,
arrondi à l'entier)} et l'affiche, suivi de « FCFA », dans la zone
\code{resultat}. \hfill \textbf{1,25pt}
\end{enumerate}

\bigskip
\textbf{Exercice 2 — Enregistrement des sessions (PHP + MySQL).} \hfill \textbf{(6 points)}

\smallskip
Le gérant dispose d'une base de données \textbf{CYBER\_DB} contenant la table
\textbf{SESSIONS} décrite ci-dessous. Un formulaire permet d'enregistrer une
nouvelle session ; un script PHP l'insère dans la base puis affiche la liste des
sessions du jour dans un tableau HTML.

\begin{center}\small
\begin{tabular}{@{}|l|l|l|@{}}
\hline
\textbf{Champ} & \textbf{Type} & \textbf{Description}\\
\hline
id        & INT (clé primaire, auto) & identifiant de la session\\
poste     & INT          & numéro de poste\\
client    & VARCHAR(50)  & nom du client\\
debut     & TIME         & heure de début\\
duree     & INT          & durée en minutes\\
montant   & INT          & montant en FCFA\\
\hline
\end{tabular}
\end{center}

\begin{enumerate}[label=\textbf{\arabic*-},leftmargin=2.2em]
\item \begin{enumerate}[label=\textbf{\alph*)},leftmargin=2em]
  \item Définir : \textbf{SGBD}, \textbf{requête SQL}, \textbf{clé primaire}. \hfill \textbf{0,75pt}
  \item Citer les \textbf{quatre} opérations de base du CRUD et donner, pour
  chacune, le mot-clé SQL correspondant. \hfill \textbf{1pt}
  \end{enumerate}
\item Écrire l'instruction \textbf{SQL} qui crée la table \textbf{SESSIONS}
décrite ci-dessus. \hfill \textbf{1pt}
\item Écrire le script \textbf{PHP} (fichier \code{enregistrer.php}) qui :
  \begin{enumerate}[label=\textbf{\alph*)},leftmargin=2em]
  \item se connecte à la base \textbf{CYBER\_DB} avec \code{mysqli} (serveur
  \code{localhost}, utilisateur \code{root}, mot de passe vide) ; \hfill \textbf{0,75pt}
  \item récupère par \code{\$\_POST} le poste, le client et la durée envoyés par
  le formulaire, calcule le montant \emph{($300 \times$ durée $\div\,60$, arrondi)}
  et \textbf{insère} la session (heure de début $=$ heure courante). \hfill \textbf{1,5pt}
  \end{enumerate}
\item Écrire la requête SQL et le code PHP qui \textbf{affichent dans un tableau
HTML} toutes les sessions enregistrées (colonnes : poste, client, durée, montant). \hfill \textbf{1pt}
\end{enumerate}

% =====================================================================
\subsection*{Partie II — Programmation procédurale en C \hfill (10 points)}

\textbf{Exercice 1 — Structure Session et chiffre d'affaires.} \hfill \textbf{(5 points)}

\smallskip
Pour le module hors-ligne du cyber, on représente une session de connexion par un
enregistrement (\code{struct}) regroupant le numéro de poste, le nom du client, la
durée (en minutes) et le montant (en FCFA). On gère un \textbf{tableau} de telles
sessions.

\begin{enumerate}[label=\textbf{\arabic*-},leftmargin=2.2em]
\item Définir : \textbf{structure} (enregistrement), \textbf{tableau de
structures}. \hfill \textbf{0,5pt}
\item Déclarer en langage~C la structure \code{Session} possédant les champs :
\code{poste} (entier), \code{client} (chaîne de 30 caractères), \code{duree}
(entier) et \code{montant} (entier). \hfill \textbf{1pt}
\item Déclarer un tableau \code{tab} de 50 sessions, puis écrire l'instruction qui
affecte au champ \code{montant} de la \textbf{première} session la valeur
\code{1500}. \hfill \textbf{0,75pt}
\item Écrire la fonction \code{int chiffreAffaires(Session t[], int n)} qui reçoit
le tableau et le nombre \code{n} de sessions, et \textbf{retourne} la somme des
montants de toutes les sessions (le chiffre d'affaires total). \hfill \textbf{1,5pt}
\item Écrire la fonction \code{int sessionPlusLongue(Session t[], int n)} qui
retourne l'\textbf{indice} de la session de plus longue durée. \hfill \textbf{1,25pt}
\end{enumerate}

\bigskip
\textbf{Exercice 2 — Lecture d'un programme C.} \hfill \textbf{(5 points)}

\smallskip
On donne ci-dessous un programme C \textbf{numéroté ligne par ligne}. Il affiche la
table de multiplication d'un nombre saisi par l'utilisateur.

\begin{codec}
1   #include <stdio.h>
2
3   int main(void) {
4       int n, i, produit;
5       printf("Entrez un nombre : ");
6       scanf("%d", &n);
7       for (i = 1; i <= 10; i++) {
8           produit = n * i;
9           printf("%d x %d = %d\n", n, i, produit);
10      }
11      return 0;
12  }
\end{codec}

\begin{enumerate}[label=\textbf{\arabic*-},leftmargin=2.2em]
\item Donner le \textbf{rôle} de chacune des lignes \textbf{4}, \textbf{6},
\textbf{7} et \textbf{8}. \hfill \textbf{1pt}
\item En supposant que l'utilisateur saisit la valeur \textbf{5}, dresser la
\textbf{trace pas-à-pas} (valeurs de \code{i} et \code{produit}) pour les
\textbf{trois premières} itérations de la boucle, puis donner la dernière ligne
affichée. \hfill \textbf{1,5pt}
\item Réécrire l'instruction de la ligne~8 en utilisant l'\textbf{opérateur
ternaire} de sorte que \code{produit} reçoive le double de \code{n*i} si \code{i}
est pair, et \code{n*i} sinon. \hfill \textbf{1pt}
\item La boucle de la ligne~7 est \textbf{itérative} (\code{for}). Écrire une
fonction \textbf{récursive} \code{void table(int n, int i)} qui affiche la même
table de multiplication (de \code{i=1} à \code{i=10}). \hfill \textbf{1,5pt}
\end{enumerate}

% =====================================================================
\corrige{du sujet TI n\textsuperscript{o}3}

\subsection*{Partie I — Programmation web}

\textbf{Ex.1 — 1) Définitions.}
\begin{itemize}
\item \emph{Page web statique} : page dont le contenu est fixe ; il est livré tel
quel au navigateur, sans traitement côté serveur ni modification dynamique.
\item \emph{Balise} : marqueur du langage HTML, écrit entre chevrons
(\code{<...>}), qui délimite et qualifie un élément de la page (ex. \code{<p>},
\code{<input>}).
\item \emph{Événement \code{onclick}} : événement déclenché par un \textbf{clic}
de l'utilisateur sur un élément ; il permet d'associer à ce clic l'appel d'une
fonction JavaScript.
\end{itemize}

\textbf{Ex.1 — 2) JavaScript vs PHP.} Un script \textbf{JavaScript} s'exécute
\textbf{côté client} (dans le navigateur du visiteur), tandis qu'un script
\textbf{PHP} s'exécute \textbf{côté serveur} (le serveur interprète le code et
n'envoie au navigateur que le HTML produit).

\textbf{Ex.1 — 3) Formulaire HTML.}

\begin{codehtml}
<!DOCTYPE html>
<html lang="fr">
<head>
  <meta charset="utf-8">
  <title>Cyber LumièreNet — Simulateur</title>
</head>
<body>
  <h1>Cyber LumièreNet</h1>
  <p>Durée (en minutes) : <input type="text" id="duree"></p>
  <button onclick="calculerCout()">Calculer</button>
  <p id="resultat"></p>

  <script src="script.js"></script>
</body>
</html>
\end{codehtml}

\textbf{Ex.1 — 4) Fonction JavaScript \code{calculerCout}.}

\begin{codejs}
function calculerCout() {
  var minutes = parseInt(document.getElementById("duree").value);
  var montant = Math.round(minutes / 60 * 300);   // 300 FCFA par heure
  document.getElementById("resultat").innerHTML = montant + " FCFA";
}
\end{codejs}

\begin{remarque}
\code{parseInt} convertit le texte saisi en entier ; \code{Math.round} arrondit le
montant. Pour 90~minutes : $90/60 \times 300 = 450$, donc « 450 FCFA ».
\end{remarque}

\medskip
\textbf{Ex.2 — 1.a) Définitions.}
\begin{itemize}
\item \emph{SGBD} (Système de Gestion de Base de Données) : logiciel qui permet de
créer, stocker, organiser et manipuler les données d'une base (ex. MySQL).
\item \emph{Requête SQL} : instruction écrite en langage SQL adressée au SGBD pour
agir sur les données (créer, lire, modifier, supprimer).
\item \emph{Clé primaire} : champ (ou ensemble de champs) qui identifie de manière
\textbf{unique} chaque enregistrement d'une table.
\end{itemize}

\textbf{Ex.2 — 1.b) Opérations CRUD et mots-clés SQL.}
\begin{center}\small
\begin{tabular}{@{}lll@{}}
\toprule
\textbf{Opération} & \textbf{Signification} & \textbf{Mot-clé SQL}\\
\midrule
Create & Créer / insérer  & \code{INSERT}\\
Read   & Lire / consulter & \code{SELECT}\\
Update & Modifier         & \code{UPDATE}\\
Delete & Supprimer        & \code{DELETE}\\
\bottomrule
\end{tabular}
\end{center}

\textbf{Ex.2 — 2) Création de la table SESSIONS.}

\begin{codesql}
CREATE TABLE SESSIONS (
    id      INT AUTO_INCREMENT PRIMARY KEY,
    poste   INT,
    client  VARCHAR(50),
    debut   TIME,
    duree   INT,
    montant INT
);
\end{codesql}

\textbf{Ex.2 — 3) Script \code{enregistrer.php} : connexion, \code{\$\_POST} et
insertion.}

\begin{codephp}
<?php
  // a) Connexion à la base CYBER_DB avec mysqli
  $cnx = new mysqli("localhost", "root", "", "CYBER_DB");
  if ($cnx->connect_error) {
      die("Échec de connexion : " . $cnx->connect_error);
  }

  // b) Récupération des données du formulaire (méthode POST)
  $poste  = $_POST["poste"];
  $client = $_POST["client"];
  $duree  = $_POST["duree"];

  // Calcul du montant : 300 FCFA / heure, arrondi
  $montant = round(300 * $duree / 60);

  // Heure de début = heure courante
  $debut = date("H:i:s");

  // Insertion de la session dans la table
  $sql = "INSERT INTO SESSIONS (poste, client, debut, duree, montant)
          VALUES ('$poste', '$client', '$debut', '$duree', '$montant')";

  if ($cnx->query($sql)) {
      echo "<p>Session enregistrée : $client, $montant FCFA.</p>";
  } else {
      echo "<p>Erreur : " . $cnx->error . "</p>";
  }
?>
\end{codephp}

\begin{remarque}
Le formulaire d'appel (fichier \code{form.html}) enverrait ces données ainsi :
\code{<form method="post" action="enregistrer.php">} avec trois champs nommés
\code{poste}, \code{client} et \code{duree}, puis un bouton de soumission.
\end{remarque}

\textbf{Ex.2 — 4) Requête de lecture et affichage dans un tableau HTML.} On
ajoute, à la suite du script précédent, la lecture de toutes les sessions :

\begin{codephp}
<?php
  // Requête de lecture
  $res = $cnx->query("SELECT poste, client, duree, montant FROM SESSIONS");

  echo "<table border='1'>";
  echo "<tr><th>Poste</th><th>Client</th><th>Durée</th><th>Montant</th></tr>";

  // Parcours des lignes et affichage
  while ($ligne = $res->fetch_assoc()) {
      echo "<tr>";
      echo "<td>" . $ligne["poste"]   . "</td>";
      echo "<td>" . $ligne["client"]  . "</td>";
      echo "<td>" . $ligne["duree"]   . " min</td>";
      echo "<td>" . $ligne["montant"] . " FCFA</td>";
      echo "</tr>";
  }
  echo "</table>";

  $cnx->close();
?>
\end{codephp}

\subsection*{Partie II — Programmation procédurale en C}

\textbf{Ex.1 — 1) Définitions.}
\begin{itemize}
\item \emph{Structure (enregistrement)} : type de données défini par le
programmeur qui regroupe sous un même nom plusieurs champs de types éventuellement
différents.
\item \emph{Tableau de structures} : tableau dont chaque case est une structure ;
il permet de stocker une collection d'enregistrements de même type (ici, plusieurs
sessions).
\end{itemize}

\textbf{Ex.1 — 2) Déclaration de la structure \code{Session}.}

\begin{codec}
typedef struct {
    int  poste;
    char client[30];
    int  duree;     // durée en minutes
    int  montant;   // montant en FCFA
} Session;
\end{codec}

\textbf{Ex.1 — 3) Tableau de 50 sessions et affectation.}

\begin{codec}
Session tab[50];
tab[0].montant = 1500;   // première session : indice 0
\end{codec}

\textbf{Ex.1 — 4) Fonction \code{chiffreAffaires}.}

\begin{codec}
int chiffreAffaires(Session t[], int n) {
    int i, total = 0;
    for (i = 0; i < n; i++) {
        total = total + t[i].montant;   // cumul des montants
    }
    return total;
}
\end{codec}

\textbf{Ex.1 — 5) Fonction \code{sessionPlusLongue}.} Elle retourne l'indice de la
session dont la durée est la plus grande.

\begin{codec}
int sessionPlusLongue(Session t[], int n) {
    int i, indiceMax = 0;
    for (i = 1; i < n; i++) {
        if (t[i].duree > t[indiceMax].duree) {
            indiceMax = i;
        }
    }
    return indiceMax;
}
\end{codec}

\begin{remarque}
On suppose $n \geq 1$. On initialise \code{indiceMax} à 0 (première session) puis
on compare chaque durée à celle de la meilleure session trouvée jusque-là.
\end{remarque}

\medskip
\textbf{Ex.2 — 1) Rôle des lignes.}
\begin{itemize}
\item \textbf{Ligne 4} : déclaration des trois variables entières \code{n}
(le nombre), \code{i} (le compteur de boucle) et \code{produit} (le résultat).
\item \textbf{Ligne 6} : lecture au clavier de la valeur de \code{n} (\code{scanf}),
rangée à l'adresse de \code{n} grâce à l'opérateur \code{\&}.
\item \textbf{Ligne 7} : début de la boucle \code{for} qui fait varier \code{i} de
1 à 10 inclus (10 itérations).
\item \textbf{Ligne 8} : calcul du produit \code{n * i} et affectation à la
variable \code{produit}.
\end{itemize}

\textbf{Ex.2 — 2) Trace pas-à-pas pour \code{n = 5}.}

\begin{center}\small
\begin{tabular}{@{}cccl@{}}
\toprule
\textbf{Itération} & \textbf{i} & \textbf{produit} & \textbf{Affichage}\\
\midrule
1 & 1 & 5  & \code{5 x 1 = 5}\\
2 & 2 & 10 & \code{5 x 2 = 10}\\
3 & 3 & 15 & \code{5 x 3 = 15}\\
\bottomrule
\end{tabular}
\end{center}

\noindent La boucle se poursuit jusqu'à \code{i = 10} ; la \textbf{dernière} ligne
affichée est \code{5 x 10 = 50}.

\textbf{Ex.2 — 3) Ligne 8 avec l'opérateur ternaire.} Le double si \code{i} est
pair, sinon \code{n*i} :

\begin{codec}
produit = (i % 2 == 0) ? 2 * n * i : n * i;
\end{codec}

\begin{remarque}
La condition \code{i \% 2 == 0} teste la parité de \code{i} : si le reste de la
division par 2 est nul, \code{i} est pair et \code{produit} reçoit \code{2*n*i} ;
sinon il reçoit \code{n*i}.
\end{remarque}

\textbf{Ex.2 — 4) Version récursive \code{table}.}

\begin{codec}
void table(int n, int i) {
    if (i > 10) {
        return;                          // cas d'arrêt
    }
    printf("%d x %d = %d\n", n, i, n * i);
    table(n, i + 1);                     // appel récursif
}
\end{codec}

\noindent On l'appelle par \code{table(n, 1);} pour obtenir la table complète de
\code{n}, identique à la version itérative.

\begin{aretenir}
Sujet en deux moitiés de 10 points : (I) un \textbf{simulateur} HTML/JavaScript
(événement \code{onclick}, lecture d'un champ, calcul, affichage) puis un module
\textbf{PHP/MySQL} (connexion \code{mysqli}, récupération par \code{\$\_POST},
\code{INSERT}, puis \code{SELECT} affiché dans un tableau HTML) ; (II) un
\textbf{tableau de \code{struct}} (\code{Session}) avec une fonction de
\textbf{chiffre d'affaires} et de \textbf{recherche du maximum}, puis la
\textbf{lecture} d'un programme C (rôles, trace, opérateur \textbf{ternaire},
passage \textbf{itératif$\rightarrow$récursif}).
\end{aretenir}
